% Ad Familiares 11.28 (ad-familiares-11-28)
% Source: https://raw.githubusercontent.com/PerseusDL/canonical-latinLit/master/data/phi0474/phi056/phi0474.phi056.perseus-lat1.xml
% Fetched by scripts/fetch_latin.py

% Dateline (Perseus): Scr. Romae paulo post ep. 27.

\ciceroLetterOpener{MATIVS CICERONI S.}

\ciceroSection{1} magnam voluptatem ex tuis litteris cepi, quod quam speraram atque optaram habere te de me opinionem cognovi ; de qua etsi non dubitabam, tamen, quia maximi aestimabam ut incorrupta maneret laborabam. conscius autem mihi eram nihil a me commissum esse quod boni cuiusquam offenderet animum. eo minus credebam plurimis atque optimis artibus ornato tibi temere quicquam persuaderi potuisse, praesertim in quem mea propensa et perpetua fuisset atque esset benevolentia. quod quoniam ut volui scio esse, respondebo criminibus, quibus tu pro me, ut par erat tua singulari bonitate et amicitia nostra, saepe restitisti.

\ciceroSection{2} nota enim mihi sunt quae in me post Caesaris mortem contulerint. vitio mihi dant quod mortem hominis necessari graviter fero atque eum quem dilexi perisse indignor ; aiunt enim patriam amicitiae praeponendam esse, proinde ac si iam vicerint obitum eius rei p. fuisse utilem. sed non agam astute ; fateor me ad istum gradum sapientiae non pervenisse ; neque enim Caesarem in dissensione civili sum secutus sed amicum ; quamquam re offendebar, tamen non deserui, neque bellum umquam civile aut etiam causam dissensionis probavi, quam etiam nascentem exstingui summe studui. itaque in victoria hominis necessari neque honoris neque pecuniae dulcedine sum captus, quibus praemiis reliqui, minus apud eum quam ego cum possent, immoderate sunt abusi. atque etiam res familiaris mea lege Caesaris de minuta est, cuius beneficio plerique, qui Caesaris morte laetantur remanserunt, in civitate. civibus victis ut parceretur aeque ac pro mea salute laboravi.

\ciceroSection{3} possum igitur qui omnis voluerim incolumis eum a quo id impetratum est perisse non indignari? cum praesertim idem homines illi et invidiae et exitio fuerint. ' plecteris ergo' inquiunt 'quoniam factum nostrum improbare audes.' O superbiam inauditam alios in facinore gloriari, aliis ne dolore quidem impunite licere! at haec etiam servis semper libera fuerunt, ut timerent, gauderent, dolerent suo potius quam alterius arbitrio ; quae nunc, ut quidem isti dictitant, 'libertatis auctores' metu nobis extorquere conantur.

\ciceroSection{4} sed nihil agunt ; nullius umquam periculi terroribus ab officio aut ab humanitate desciscam ; numquam enim honestam mortem fugiendam, saepe etiam oppetendam putavi. sed quid mihi suscensent, si id opto ut paeniteat eos sui facti? cupio enim Caesaris mortem omnibus esse acerbam. ' at debeo pro civili parte rem p. velle salvam.' id quidem me cupere, nisi et ante acta vita et reliqua mea spes tacente me probat, dicendo vincere non postulo.

\ciceroSection{5} qua re maiorem in modum te rogo ut rem potiorem oratione ducas mihique, si sentis expedire recte fieri, credas nullam communionem cum improbis esse posse. an quod adulescens praestiti, cum etiam errare cum excusatione possem, id nunc aetate praecipitata commutem ac me ipse retexam? non faciam neque quod displiceat committam, praeterquam quod hominis mihi coniunctissimi ac viri amplissimi doleo gravem casum. quod si aliter essem animatus, numquam quod facerem negarem, ne et in peccando improbus et in dissimulando timidus ac vanus existimarer.

\ciceroSection{6} ' at ludos quos Caesaris victoriae Caesar adulescens fecit curavi.' at id ad privatum officium, non ad statum rei p. pertinet ; quod tamen munus et hominis amicissimi memoriae atque honoribus praestare etiam mortui debui et optimae spei adulescenti ac dignissimo Caesare petenti negare non potui.

\ciceroSection{7} veni etiam consulis Antoni domum saepe salutandi causa ; ad quem qui me parum patriae amantem esse existimant rogandi quidem aliquid aut auferendi causa frequentis ventitare reperies. sed quae haec est adrogantia, quod Caesar numquam interpellavit quin quibus vellem atque etiam quos ipse non diligebat tamen iis uterer, eos qui mihi amicum eripuerunt carpendo me efficere conari ne quos velim diligam?

\ciceroSection{8} sed non vereor ne aut meae vitae modestia parum valitura sit in posterum contra falsos rumores aut ne etiam ii, qui me non amant propter meam in Caesarem constantiam, non malint mei quam.sui similis amicos habere. mihi quidem si optata contingent, quod reliquum est vitae in otio Rhodi degam ; sin casus aliquis interpellarit, ita ero Romae ut recte fieri semper cupiam. Trebatio nostro magnas ago gratias, quod tuum erga me animum simplicem atque amicum aperuit et quod eum, quem semper libenter dilexi, quo magis iure colere atque observare deberem fecit. bene vale et me dilige.
